(1) 在小时间 \(dt\) 内火箭质量降低 \(dm<0\)，排出气体在惯性系的向上速度近似 \(v-u\)。对火箭与气体组成的瞬时系统作竖直动量平衡并保留一阶量，有 \(m\,dv=-u\,dm-mg\,dt\)，即 \(\dot v=u\dot m/m-g\)（这里 \(\dot m>0\) 表示正的燃料排出速率，实际 \(dm/dt=-\dot m\)）。
(2) 由 \(m(t)=m_0-\dot m t\)，燃尽时间 \(t_b=(m_0-m_f)/\dot m\)。对 \(\dot v=u\dot m/(m_0-\dot m t)-g\) 积分，得到 \(v(t)=u\ln[m_0/m(t)]-gt\)，从而燃尽速度 \(v_b=u\ln(m_0/m_f)-g(m_0-m_f)/\dot m\)。后一项正是重力损失而非排气冲量。
(3) 积分 \(z(t)=\int_0^t v(s)ds\)。利用 \(\int_0^t\ln[m_0/(m_0-\dot m s)]ds=t-[m(t)/\dot m]\ln[m_0/m(t)]\)，得 \(z(t)=u\{t-[m(t)/\dot m]\ln[m_0/m(t)]\}-gt^2/2\)。故 \(z_b=u\{t_b-(m_f/\dot m)\ln(m_0/m_f)\}-gt_b^2/2\)。
(4) 燃尽后速度按 \(v(t_b+\tau)=v_b-g\tau\) 下降。若 \(v_b>0\)，到最高点的额外上升高度为 \(\Delta z=v_b^2/(2g)\)；若 \(v_b\le0\)，则火箭在燃尽时不再上升，不能把该式误作正的未来爬升。\(g\to0\) 时燃烧速度恢复 \(u\ln(m_0/m_f)\)，但燃尽后没有有限最高点；\(u\to0\) 时仅剩重力自由落体。
